\section*{Chapitre \ref{ch:g4:measure} --- Mesures et périmètre}

\begin{solution}{exo:g4:measure:1}
$7$~m $= 700$~cm~; \quad $3.2$~km $= 3\,200$~m~; \quad
$85$~mm $= 8.5$~cm~; \quad $4\,500$~m $= 4.5$~km.
\end{solution}

\begin{solution}{exo:g4:measure:2}
$2$~kg $= 2\,000$~g~; \quad $1\,800$~g $= 1$~kg $800$~g~; \quad
$3.5$~L $= 350$~cL~; \quad $25$~cL $= 0.25$~L.
\end{solution}

\begin{solution}{exo:g4:measure:3}
Rectangle~: $2 \times (9 + 5) = 28$~cm. Carré~:
$4 \times 7.5 = 30$~cm. Triangle~: $6 + 8 + 10 = 24$~cm.
\end{solution}

\begin{solution}{exo:g4:measure:4}
Longueur $+$ largeur $= 26 \div 2 = 13$~cm, donc la largeur est
$13 - 8 = 5$~cm.
\end{solution}

\begin{solution}{exo:g4:measure:5}
De nombreuses formes en L conviennent~; pour un L de $10$ carreaux
fait d'un rectangle $4 \times 2$ plus un pied $2 \times 1$,
parcourir le bord donne $16$ unités. (Toute figure correcte et un
comptage soigneux sont justes~; le périmètre dépend du L choisi.)
\end{solution}

\begin{solution}{exo:g4:measure:6}
$3 \times 4$~: périmètre $14$. $2 \times 6$~: périmètre $16$.
($1 \times 12$~: périmètre $26$.) Plus on se rapproche d'un carré,
plus le périmètre est petit pour une même aire.
\end{solution}

\begin{solution}{exo:g4:measure:7}
Moitié du rectangle $6 \times 3$~: $18 \div 2 = 9$ carreaux.
\end{solution}

\begin{solution}{exo:g4:measure:8}
$10{:}20 \to 12{:}45$~: $2$~h $25$~min. $7{:}35 \to 9{:}10$~:
$1$~h $35$~min. Film~: $20{:}30 + 1$~h $50 = 22{:}20$.
\end{solution}

\begin{solution}{exo:g4:measure:9}
Un tour~: $2 \times (60 + 45) = 210$~m. Trois tours~:
$3 \times 210 = 630$~m.
\end{solution}

\begin{solution}{exo:g4:measure:10}
Périmètre~: $4 \times 85 = 340$~m~; moins le portail~:
$340 - 4 = 336$~m de clôture.
\end{solution}

\begin{solution}{exo:g4:measure:11}
Les deux affirmations sont \emph{fausses}. Même périmètre, aires
différentes~: le rectangle $4 \times 2$ et la forme en L de la
figure du chapitre (tous deux périmètre $12$~; aires $8$ et $6$).
Même aire, périmètres différents~: le carré $4 \times 4$ et le
rectangle $8 \times 2$ (tous deux aire $16$~; périmètres $16$ et
$20$).
\end{solution}
